%% OUP authoring template — Bioinformatics Application Note
%% Requires: oup-authoring-template.cls (provided by the Overleaf template)
%% Template: https://www.overleaf.com/latex/templates/oup-general-template/ybpypwncdxyb

\documentclass[unnumsec,webpdf,contemporary,large]{oup-authoring-template}

%%% ---------------------------------------------------------------
%%% Packages
%%% ---------------------------------------------------------------
\usepackage{microtype}
\usepackage{hyperref}
\usepackage{xurl}          % allow line-breaks in long URLs
\usepackage{booktabs}
\usepackage{graphicx}

% \authormark is referenced by the class internally but never defined
\newcommand{\authormark}[1]{\markboth{#1}{#1}}
% \city is missing from the class address helpers
\newcommand{\city}[1]{#1}

%%% ---------------------------------------------------------------
%%% Journal metadata (fill in after acceptance)
%%% ---------------------------------------------------------------
\journaltitle{Bioinformatics}
\copyrightyear{2026}
\pubyear{2026}
\appnotes{Application Note}

%%% ---------------------------------------------------------------
%%% Title
%%% ---------------------------------------------------------------
\title[interactivePCA]{interactivePCA: simultaneous exploration of
  ancient DNA across genetic, geographic, temporal, and phenotypic
  dimensions}

%%% ---------------------------------------------------------------
%%% Authors and affiliations
%%% ---------------------------------------------------------------
\author[1,2,*]{Samuel Neuenschwander}
\author[1,3]{Anna-Sapfo Malaspinas}

\authormark{Neuenschwander and Malaspinas}

\address[1]{%
  \orgdiv{Department of Computational Biology},
  \orgname{University of Lausanne},
  \orgaddress{\city{Lausanne}, \country{Switzerland}}}
\address[2]{%
  \orgdiv{Vital-IT},
  \orgname{SIB Swiss Institute of Bioinformatics},
  \orgaddress{\city{Lausanne}, \country{Switzerland}}}
\address[3]{%
  \orgname{SIB Swiss Institute of Bioinformatics},
  \orgaddress{\city{Lausanne}, \country{Switzerland}}}

\corresp[*]{%
  Corresponding author.
  \href{mailto:samuel.neuenschwander@unil.ch}%
       {samuel.neuenschwander@unil.ch}}

%%% ---------------------------------------------------------------
%%% Abstract and keywords
%%% ---------------------------------------------------------------
\abstract{%
Principal component analysis (PCA) is the central exploratory tool in
ancient DNA (aDNA) research, yet standard implementations constrain
analysts to a single two-dimensional view at a time.
We present \textit{interactivePCA}, a browser-based application that
links up to five analytical dimensions in a single coordinated
interface: two genetic principal components, geographic provenance,
chronological age, and a user-defined grouping variable.
Lasso selections, hover information, and aesthetic styling propagate
instantly across all panels, enabling cohesive visual interrogation of
population structure, migration, and admixture.
\textit{interactivePCA} is implemented in Python/Dash, requires no
server infrastructure, and accepts the PLINK eigenvector format
directly.}

\keywords{ancient DNA, principal component analysis, interactive
  visualisation, population genetics, multidimensional exploration}

%%% ---------------------------------------------------------------
%%% Availability
%%% ---------------------------------------------------------------
\otherabstract[Availability and implementation]{%
Source code and documentation are freely available at
\url{https://github.com/sneuensc/interactive-pca} under the MIT
licence. The tool requires Python~$\ge 3.8$, Dash~$\ge 2.0$,
Plotly~$\ge 5.0$, pandas~$\ge 1.3$, and dash-ag-grid~$\ge 2.0$.}

%%%================================================================
\begin{document}
\maketitle
%%%================================================================

\section{Introduction}

Ancient DNA studies have transformed our understanding of human
prehistory by revealing large-scale migrations, admixture events, and
population replacements that left no unambiguous trace in the
archaeological record
\citep{Haak2015,Mathieson2015,Reich2018}.
PCA remains the workhorse for presenting genetic variation: it is
computationally tractable on cohorts of thousands of individuals and
produces intuitive scatter plots that mirror population structure
\citep{Patterson2006}.

However, aDNA data are inherently multi-dimensional.
A single ancient individual carries simultaneously a genomic signature
(its position in PC space), a geographic provenance (site
latitude/longitude), a chronological age (radiocarbon date or
archaeological period), and one or more categorical descriptors
(culture, haplogroup, sex).
Standard PCA software---EIGENSOFT \citep{Price2006} and PLINK
\citep{Purcell2007}---writes static eigenvector tables; visualisation
is then delegated to general-purpose plotting packages such as
R/ggplot2 or Python/matplotlib.
Static images, however, force analysts to cross-reference separate
figures mentally, a process that is slow, error-prone, and becomes
impractical when cohort sizes reach hundreds of samples.

Interactive web-based tools such as HGDP Explorer and GenomePlot
provide some dynamic functionality, but none natively links a PC
scatter plot with a geographic map and a time axis in a fully
coordinated, lasso-selectable interface suitable for local, offline
use with private data.
\textit{interactivePCA} fills this gap, allowing researchers to select
a subset of individuals in any one panel and immediately see exactly
where those individuals lived, when they lived, and how they cluster in
PC space---simultaneously and without writing a single line of code.

%%-----------------------------------------------------------------
\section{Implementation}
%%-----------------------------------------------------------------

\textit{interactivePCA} is a Python package built on the Dash reactive
web framework \citep{Dash2024} with Plotly.js for rendering.
All computation is performed at startup; once the application is
running, every interaction---selections, hover events, axis
changes---is handled client-side through compiled JavaScript callbacks,
giving sub-100\,ms response times even on datasets of several thousand
samples.

\paragraph{Input.}
The tool accepts two files: (i)~a PLINK-format eigenvector file
(\texttt{--eigenvec}) containing any number of principal components,
and (ii)~an optional tab-separated annotation file
(\texttt{--annotation}) whose column roles---identifier, latitude,
longitude, chronological date, grouping variables---are specified via
command-line flags or auto-detected by column name.
Pre-computed MDS coordinates are equally accepted.

\paragraph{Architecture.}
The application is launched via a single CLI command
(\texttt{interactive-pca --eigenvec \ldots{} --annotation \ldots{}})
and served on \texttt{localhost}, ensuring that confidential genomic
data never leave the analyst's machine.
The layout consists of two resizable panes: the left pane holds the
PCA/MDS scatter plot and the time scatter plot (vertically stacked,
with a draggable divider); the right pane holds the geographic map and
a tabbed panel containing the annotation table, a hover-detail view,
and a pandas query filter.
A clientside JavaScript callback propagates hover and selection events
across all three plots in real time without server round-trips.

\paragraph{Customisation.}
Point colours, symbols, sizes, and opacities are configurable per
group value through a modal aesthetics editor and can be persisted to a
JSON file for reproducibility.
Continuous variables (e.g.\ radiocarbon dates, allele frequencies) are
rendered with user-selectable colour scales; categorical grouping
variables receive a default qualitative palette that the user can
override.

%%-----------------------------------------------------------------
\section{Usage and key features}
%%-----------------------------------------------------------------

\subsection{Five simultaneous analytical dimensions}

A typical aDNA session opens with a two-dimensional PC scatter plot
(dimensions 1 and 2), a Mapbox-backed geographic scatter plot showing
sample provenance, and a chronological scatter plot (jittered along
the $y$-axis to separate overlapping dates).
Taken together, these three panels expose \textbf{five analytical
dimensions at once}: PC1 ($x$), PC2 ($y$), longitude, latitude, and
time.
A sixth, discrete dimension is encoded through the \textit{Group~by}
selector, which re-colours all three panels simultaneously according to
any categorical or continuous annotation column---culture, haplogroup,
admixture proportion, or quality metric.

\subsection{Coordinated cross-panel selection}

Lasso or box selection in any panel instantly highlights the matched
individuals in the other two panels and updates the annotation table.
Conversely, toggling checkboxes in the annotation table, or applying a
pandas query filter (e.g.\ \texttt{Date < -3000 \& Culture ==
"Yamnaya"}), propagates the selection back to all three plots.
This bidirectionality enables rapid hypothesis testing: an analyst can,
for example, isolate genomically outlying individuals from a PC cluster
and immediately inspect whether they share a geographic or
chronological signal.

\subsection{Hover information and detail panel}

Hovering over any point in any plot simultaneously displays the
matching point in the other two plots, showing at minimum the sample
identifier and group label.
A \textit{Hover detailed} toggle enriches the tooltip with all
annotation columns currently selected in the annotation table.
When the side-panel \textit{Details} tab is active, the hover card is
expanded into a scrollable key--value table, and the in-plot tooltips
revert to minimal labels to avoid visual clutter.

\subsection{Three-dimensional PCA and axis switching}

PC axes are freely switchable via drop-down menus ($x$, $y$, and, in
3D mode, $z$), allowing users to traverse the full PC space without
regenerating figures.
Three-dimensional mode renders an interactive Scatter3d widget in which
the lasso tool is replaced by 3D camera rotation; selections made in
3D are propagated to the 2D geographic and time panels.

%%-----------------------------------------------------------------
\section{Conclusion}
%%-----------------------------------------------------------------

\textit{interactivePCA} provides an integrated, five-dimensional visual
workspace for aDNA research that substantially reduces the cognitive
overhead of multi-panel analysis.
By linking genetic, geographic, and chronological dimensions within a
single coordinated interface, it enables analysts to detect
patterns---such as the co-occurrence of PC outliers with
chronologically early or geographically peripheral
samples---that would be invisible in separate static plots.
The tool is designed to integrate directly into existing
PLINK/EIGENSOFT workflows, requires no bioinformatics infrastructure
beyond a standard Python environment, and scales to the cohort sizes
typical of current aDNA studies (hundreds to a few thousand samples).

%%-----------------------------------------------------------------
\section*{Acknowledgements}

The authors thank members of the Malaspinas lab for testing and
feedback.

\section*{Funding}

[Funding sources and grant numbers to be added.]

%%%================================================================
%%% Bibliography
%%%================================================================
\begin{thebibliography}{9}

\bibitem[Dash(2024)]{Dash2024}
Plotly Technologies Inc. (2024)
\newblock Dash: analytical web apps for Python, R, Julia, and Jupyter.
\newblock \url{https://dash.plotly.com}

\bibitem[Haak et al.(2015)]{Haak2015}
Haak,~W. et al. (2015)
\newblock Massive migration from the steppe was a source for Indo-European
  languages in Europe.
\newblock \textit{Nature}, \textbf{522}, 207--211.

\bibitem[Mathieson et al.(2015)]{Mathieson2015}
Mathieson,~I. et al. (2015)
\newblock Genome-wide patterns of selection in 230 ancient Eurasians.
\newblock \textit{Nature}, \textbf{528}, 499--503.

\bibitem[Patterson et al.(2006)]{Patterson2006}
Patterson,~N. et al. (2006)
\newblock Population structure and eigenanalysis.
\newblock \textit{PLOS Genetics}, \textbf{2}, e190.

\bibitem[Price et al.(2006)]{Price2006}
Price,~A.L. et al. (2006)
\newblock Principal components analysis corrects for stratification in
  genome-wide association studies.
\newblock \textit{Nature Genetics}, \textbf{38}, 904--909.

\bibitem[Purcell et al.(2007)]{Purcell2007}
Purcell,~S. et al. (2007)
\newblock PLINK: a tool set for whole-genome association and population-based
  linkage analyses.
\newblock \textit{American Journal of Human Genetics}, \textbf{81}, 559--575.

\bibitem[Reich(2018)]{Reich2018}
Reich,~D. (2018)
\newblock \textit{Who We Are and How We Got Here: Ancient DNA and the New
  Science of the Human Past}.
\newblock Oxford University Press, Oxford.

\end{thebibliography}

\end{document}
